\section{A Transport That Works: The Bit-Banged Master}
\label{sec:i2c-bitbang}

Once the on-chip controller is out of the picture, the bus becomes software, and the
software can be made small enough to hold in your head.

\subsection{The Whole Protocol in Four Primitives}

Every I2C transfer — a register write, a 12-byte burst, a scan of 128 addresses — reduces
to four operations on the two lines:

\begin{minted}{rust}
/// Let SDA float (the pull-up makes it a 1) or pull it to ground.
fn sda(&mut self, high: bool);
/// Same for the clock.
fn scl(&mut self, high: bool);
/// Let the bus settle for half a bit period, then sample SDA back.
fn read_sda(&self) -> bool;
/// Half a bit period of wall time.
fn half_bit(&self);
\end{minted}

With those, a byte is a loop and a START is two writes:

\begin{minted}{rust}
fn start(&mut self) -> bool {
    self.sda(true);
    self.scl(true);
    self.half_bit();
    self.sda(false);        // SDA falls while SCL is high: START
    self.half_bit();
    self.scl(false);
    self.half_bit();
    true
}

fn write_byte(&mut self, byte: u8) -> bool {
    for bit in (0..8).rev() {
        self.sda(byte & (1 << bit) != 0);
        self.half_bit();
        self.scl(true);
        self.half_bit();
        self.scl(false);    // data is sampled by the device on the rising edge
    }
    // Ninth clock: release SDA and see whether anybody pulls it down.
    self.sda(true);
    self.half_bit();
    self.scl(true);
    self.half_bit();
    let acked = !self.read_sda();
    self.scl(false);
    acked
}
\end{minted}

\begin{tipbox}
    The acknowledge is the only place where the master \emph{reads}. Reading SDA during a
    write is what allows the driver to report ``nobody is home'' instead of silently
    succeeding — and that distinction is what makes a bus scan possible at all.
\end{tipbox}

\subsection{Timing Without a Timer}

The half-bit period does not need an interrupt or a timer peripheral. A counted loop is
enough, and its length is calibrated against something known: the CPU clock.

\begin{minted}{rust}
/// ~5 us at the 19 MHz this board runs at, which is a 100 kHz bus.
const HALF_BIT_CYCLES: u32 = 60;

fn half_bit(&self) {
    cortex_m::asm::delay(HALF_BIT_CYCLES);
}
\end{minted}

\begin{warningbox}
    \codeinline{cortex\_m::asm::delay(n)} does \textbf{not} execute exactly $n$ CPU cycles.
    On this part it ran about three times faster than the name suggests, which turned a
    nominal 100 kHz bus into a 300 kHz one — and turned a buzzer calibration into a 3$\times$
    error elsewhere in the same project. Anything that claims a frequency should be measured
    with the cycle counter (\codeinline{DWT->CYCCNT}), not assumed from a loop count.
\end{warningbox}

\subsection{The \texttt{embedded-hal} Face}

The transport is more useful if the rest of the ecosystem can use it. \codeinline{embedded-hal}
v1.0 defines the trait, and the implementation is a thin wrapper over the primitives above:

\begin{minted}{rust}
impl<B: I2cBus> embedded_hal::i2c::I2c for SharedBus<'_, B> {
    fn transaction(&mut self, address: u8, operations: &mut [Operation<'_>])
        -> Result<(), Error>
    {
        // One START, the operations in order, one STOP. A register read is a write of the
        // register number followed by a repeated START and a read: getting the repeated
        // START wrong is the second most common cause of a silent bus.
        ...
    }
}
\end{minted}

Designing the driver against the trait rather than against the pin-level functions is what
allowed the same SSD1306, BME68x, PCF8563 and LSM6DS3TR-C drivers to be written once and
used on two different buses — one on PC04/PC05, another on PB2/PB3 — without changing a line
of driver code.

\subsection{Two Buses on One Board}

The board this book was written against has two I2C buses, and the reason is worth recording:

\begin{table}[H]
\centering
\begin{tabular}{llll}
\hline
\textbf{Bus} & \textbf{Pins} & \textbf{Devices} & \textbf{Notes} \\
\hline
I2C on the base & PC04 (SDA), PC05 (SCL) & OLED 0x3C, BME68x 0x76, RTC 0x51 & pull-ups on the base \\
IMU bus         & PB2 (SDA), PB3 (SCL)  & LSM6DS3TR-C 0x6A                & external pull-ups, own power gate \\
\hline
\end{tabular}
\caption{The two I2C buses of the XIAO MG24 Sense on the XIAO Expansion Base.}
\label{tab:i2c-two-buses}
\end{table}

That second row is where the chapter's hardest lesson lives, and it is the subject of the
next section.
